% !TEX root = ../main.tex

\section{Method}
\label{sec:method}

\begin{figure*}[t]
\centering
\includegraphics[width=0.95\textwidth]{figures/framework.pdf}
\caption{Overview of the Adaptive Auto-Harness system. \textbf{Top:} the multi-agent evolver constructs and refines a harness tree across cycles via four phases (Analyze $\to$ Research $\to$ Build $\to$ Verify) with a persistent cross-cycle workspace and temporal-reveal feedback. \textbf{Bottom:} at solve time, a router agent reads each branch's workspace via \texttt{git show} and routes the incoming task $x_t$ to the most suitable branch. Two human-in-the-loop hooks (task-board steering and research-phase assistance) trigger only when the evolver's history lacks relevant signal.}
\label{fig:unified_pipeline}
\end{figure*}


\subsection{Problem Formulation}
\label{sec:method-setup}

We consider tasks arriving in an open-ended stream $x_1, x_2, \ldots, x_T$ with $x_t \sim P_t$, where $P_t$ is the task distribution at time $t$ and each $x_t$ has a fixed ground truth $y(x_t)$. The agent observes the history of prior experience $\Hist = \{(x_i, r_i, \tau_i)\}_{i=1}^{t-1}$, where $r_i$ is optionally the realized reward on task $x_i$ and $\tau_i$ is the agent's solving trajectory (actions, intermediate observations, tool calls). A \emph{harness} $C = \varphi(\Hist)$, with $|C| \leq K$, is a bounded representation (prompts, skills, memory, and tools) with capacity budget $K$, where $\varphi$ is the \emph{evolver}: the agentic system that automatically transforms historic experience into the harness. A solver agent then acts via the policy $a_t \sim \pi(a \mid x_t, C_t)$, where $\pi$ is the LLM-induced sampling distribution over actions conditioned on the task and harness.

\subsection{Analytical Framework}
\label{sec:framework}

The three deployment dimensions of section \S\ref{sec:intro} surface concrete failures, but they do not by themselves identify what an evolver should fix. To pinpoint the root causes, we frame the problem analytically: we cast harness construction as a regret-minimization problem against an oracle reference, and decompose the regret into two complementary loss terms that map directly to actionable design axes.

\noindent\textbf{Utility and regret.}
We define the utility of a harness $C$ on a task $x_t$ as
\begin{equation}
\label{eq:utility}
V(C, x_t) = \mathbb{E}_{a_t \sim \pi(\cdot \mid x_t, C)}[r(a_t, y(x_t))],
\end{equation}
which is the expected reward of the harness-conditioned solver. The full-history utility is the corresponding ceiling under the capacity budget,
$V(\Hist, x_t) \,:=\, \sup_{\substack{C \,:\, |C| \leq K}} V(C, x_t)$,
that is, the best utility attainable by any bounded harness on $x_t$. The regret of the evolver-constructed harness is then
\begin{equation}
\label{eq:regret}
\text{Regret}(\varphi, x_t) = V(\Hist, x_t) - V(\varphi(\Hist), x_t) \geq 0,
\end{equation}
non-negative since $\varphi(\Hist)$ is one element of the supremum's domain. Operationally, this ceiling matches what a solver granted $\Hist$ directly, under the same compute and tool budget as $\varphi$, could attain by reconstructing any candidate harness on the fly.

\noindent\textbf{Solve-time optimal harness.}
Fix the evolver class $\Phi$ and define the solve-time optimal harness for a particular task:
\begin{equation}
\label{eq:cstar}
C^*_\Phi(x_t) = \arg\max_{\substack{C = \varphi(\Hist, x_t),\ \varphi \in \Phi \\ |C| \leq K}} V(C, x_t).
\end{equation}
This is an hypothetic oracle reference: it is the best harness the evolver class $\Phi$ could produce if it were allowed to condition on the incoming task $x_t$ at solve time. A deployed evolver $\varphi$ commits to one harness $\varphi(\Hist)$ before $x_t$ is observed; using $C^*_\Phi$ as pivot between $V(\Hist, x_t)$ and $V(\varphi(\Hist), x_t)$ yields the following decomposition.

\begin{proposition}[Regret Decomposition]
\label{prop:decomp}
For any deployed evolver $\varphi \in \Phi$,
\begin{equation}
\label{eq:decomp}
\mathbb{E}_{x_t}[\text{Regret}(\varphi, x_t)] = \Levo(\Phi) + \Ladapt(\varphi),
\end{equation}
where
\begin{align}
\label{eq:losses}
\Levo(\Phi) &= \mathbb{E}_{x_t}\!\bigl[V(\Hist, x_t) - V(C^*_\Phi(x_t), x_t)\bigr],
\end{align}
\begin{align}
\Ladapt(\varphi)  &= \mathbb{E}_{x_t}\!\bigl[V(C^*_\Phi(x_t), x_t) - V(\varphi(\Hist), x_t)\bigr].
\end{align}
\end{proposition}

$\Levo$ is the \emph{evolution loss}: it reflects the evolver class $\Phi$'s capability gap, regardless of how many cycles the evolver runs. A single-agent prompt editor cannot produce multi-file infrastructure; that ceiling is structural, not a matter of effort. Reducing $\Levo$ therefore requires pursuing more capable evolver systems with access to broader control and feedbacks. $\Ladapt$ is the \emph{adaptation loss}: the oracle builds the optimal harness per task, but a deployed $\varphi$ commits to one harness before seeing $x_t$. Therefore, even with an optimal evolver system, $\Ladapt$ exists as long as task heterogeneity persists.

% Since $C^*_\Phi(x_t)$ dominates any fixed harness pointwise,
% \begin{equation}
% \label{eq:conditioning_gap}
% \mathbb{E}_{x_t}\!\bigl[V(C^*_\Phi(x_t), x_t)\bigr] \;\geq\; \max_{\varphi \in \Phi}\, \mathbb{E}_{x_t}\!\bigl[V(\varphi(\Hist), x_t)\bigr],
% \end{equation}
% with strict inequality whenever different tasks prefer different harnesses. 

\noindent\textbf{Human-in-the-loop as a third axis.}
The decomposition assumes $\Hist$ contains relevant signal for $x_t$. When it does not, the regret decomposition no longer applies; we address this case via a third axis outside $\Levo + \Ladapt$, the human-in-the-loop channel (\S\ref{sec:method-hitl}).


\subsection{Sustained Auto-Harness via Multi-Agent Evolution}
\label{sec:method-multi}

The analytical decomposition identifies $\Levo$ as the loss from harness capabilities that the evolver class cannot construct from history. We reduce this loss by expanding $\Phi$ with a stateful four-phase multi-agent evolver, temporal-reveal feedback, and cross-cycle memory (Figure~\ref{fig:unified_pipeline}). This design targets the unbounded-stream failure where a single-agent evolver must absorb growing trajectories, delayed labels, and prior research within one context window. Concretely, we address three structural limitations:

\noindent\textbf{(1) Multi-agent with distinct roles and objectives.}
Because unbounded streams keep expanding the trajectory history an evolver must interpret, existing single-agent evolvers must fit analysis, research, implementation, and verification into one context window. We decompose evolution into four phases (Analyze $\to$ Research $\to$ Build $\to$ Verify), each with a dedicated objective and full context budget. This eliminates the single-window bottleneck and lets parallel Researchers explore disjoint hypotheses without premature convergence.

\noindent\textbf{(2) Temporal-reveal feedback.}
Under unbounded streams, labels arrive asynchronously (e.g., a prediction market resolves days after the trade). We implement a \emph{temporal-reveal gate} that surfaces each task's evaluation signal to the evolver only after its resolution date, providing a proper streaming feedback signal without leaking future information.

\noindent\textbf{(3) Persistent cross-cycle state.}
We provide the evolver with a dedicated workspace that persists across cycles, containing a task board (prioritised failure analysis), research logs (tested hypotheses with pass/fail verdicts), architecture documentation (README), and verification tests. This cross-cycle memory enables the evolver to refine its construction ability over time and build upon prior evolution experience rather than restarting from scratch.




\subsection{Solve-Time Adaptation via Harness-Tree Routing}
\label{sec:method-nav}

The analytical decomposition identifies $\Ladapt$ as the loss from committing to one dense harness before observing the incoming task's context. We reduce this loss by shifting heavy adaptation into evolution time: the evolver constructs a structured harness store, and solve time only requires a lightweight adaptation operator. The reduction is general: it factors into (i) how the evolver organizes the harness space and (ii) how the solver selects from that space per task. The harness space can be organized in many forms, including a linear chain that always uses the most recently evolved workspace, a tree of regime-specific branches, or a graph of skills with dependency edges. The adaptation operator can likewise take many forms, ranging from skill-level retrieval over a flat catalog~\cite{deerflow2025} to branch-level routing over a structured space. Different combinations trade off construction cost, operator latency, and the granularity at which adaptation occurs.

We adopt a \emph{harness tree} as the storage form and \emph{agentic routing} as the adaptation operator. The tree is the natural fit for our setting because heterogeneous task streams cluster into a small number of recurring regimes (e.g., binary exploitation versus cryptography in CTF-Dojo, sports versus politics in PolyBench), branches isolate regime-specific prompts/skills/tools without cross-contaminating the others, and the branching gate gives the evolver an explicit lever to commit specialization only when warranted by failure evidence. We instantiate this with two designs (Figure~\ref{fig:unified_pipeline}). \textbf{(1) Branching harness tree (built at evolution time).} The solver's workspace is a git repository. The evolver constructs regime-specific branches (e.g., \texttt{branch/crypto-classical}, \texttt{branch/binary-reversing}) during evolution, each carrying its own prompt, skills, and tool registry; git provides versioning, isolation, and lineage tracking across branches. \textbf{(2) Agentic routing (executed at solve time).} A router agent reads each branch's workspace via \texttt{git show} and selects the branch given $x_t$'s context; the solver then checks out that branch and executes.

\subsection{Human-in-the-Loop Channel}
\label{sec:method-hitl}

A third failure lies outside $\Levo + \Ladapt$: some tasks require harness absent from $\Hist$ signal, such as API credentials, novel web sources, or proprietary endpoints. In this experience-insufficient setting, neither a stronger evolver nor solve-time routing can recover the missing signal. We address it with a human-in-the-loop channel that augments $\Hist$ through structurally triggered steering hooks (Figure~\ref{fig:unified_pipeline}). This design targets open-ended streams where new access requirements appear before autonomous evolution has relevant evidence.
\noindent\textbf{(1) Task-board steering.}
After the Analyst updates the task board, a human may review it to add entries, adjust priorities, or supply domain guidance and source access. This proactively steers the subsequent research cycle with direction the evolver cannot derive from trajectories alone.
\noindent\textbf{(2) Interactive assistance during research.}
When a Researcher agent hits a barrier mid-execution that requires human intervention (e.g., an authentication wall), the hook prompts the human in real time. This reactively unblocks the research agent at the point of failure.
